\section{Clause-Dual Distance Wall and Conditional EntropyDepth Lower Bound}

\subsection{Linear-algebraic model}

Let $A \in \mathbb{F}_2^{m \times n}$ be a fixed matrix and write $\langle u,v\rangle = u^\top v \in \mathbb{F}_2$.
For $w \in \mathbb{F}_2^m$, write $\|w\|_0 := |\{i \in [m] : w_i = 1\}|$.

\begin{definition}[Clause-dual code and minimum clause-dual distance]
Define the \emph{clause-dual code} by
\[
\mathcal{C}^\perp_{\mathrm{cl}}(A) := \ker(A^\top) \subseteq \mathbb{F}_2^m.
\]
Define the \emph{minimum clause-dual distance} by
\[
d^\perp_{\mathrm{cl}}(A) := \min\{\|w\|_0 : w \in \mathcal{C}^\perp_{\mathrm{cl}}(A)\setminus\{0\}\}.
\]
\end{definition}

\subsection{Parity-conditioned obstruction on clause-space}

Fix a nonzero $w \in \mathcal{C}^\perp_{\mathrm{cl}}(A)$ and set $s := \|w\|_0$.
Let $\mu_0,\mu_1$ be the uniform distributions on $\mathbb{F}_2^m$ conditioned on $\langle w,b\rangle = 0$ and $\langle w,b\rangle = 1$, respectively.

\begin{proposition}[Exact local indistinguishability under a clause-parity cut]
Let $T \subseteq [m]$ with $|T| \le s-1$.
For any Boolean test $f:\mathbb{F}_2^T \to \{0,1\}$, viewing $f$ as a function on $\mathbb{F}_2^m$ via restriction to coordinates in $T$, one has
\[
\mathbb{P}_{b\sim \mu_0}[f(b_T)=1]
=
\mathbb{P}_{b\sim \mu_1}[f(b_T)=1].
\]
Equivalently, every $T$-measurable Boolean test has parity-cut bias exactly zero.
Moreover, $\|\mu_0-\mu_1\|_{\mathrm{TV}}=1$.
\end{proposition}

\begin{proof}
Let $U := \mathrm{supp}(w)$ so $|U|=s$.
Fix any $a \in \mathbb{F}_2^T$ and count completions $b \in \mathbb{F}_2^m$ with $b_T=a$.
Since $|T|\le s-1$, one has $U\not\subseteq T$, so choose a free coordinate in $U\setminus T$.
Flipping that coordinate is a bijection between completions with $\langle w,b\rangle=0$ and completions with $\langle w,b\rangle=1$ while leaving $b_T=a$ unchanged.
Hence
\[
\mu_0[b_T=a]=\mu_1[b_T=a]
\]
for every $a\in\mathbb F_2^T$.
Summing over the level set $f^{-1}(1)$ gives the displayed equality.
Finally, $\mu_0$ and $\mu_1$ have disjoint supports since $\langle w,b\rangle$ is constant on each, hence $\|\mu_0-\mu_1\|_{\mathrm{TV}}=1$.
\end{proof}

\begin{corollary}[Clause-dual distance wall for bounded-support clause tests]
Let $A\in \mathbb{F}_2^{m\times n}$ and let $d := d^\perp_{\mathrm{cl}}(A)$.
For any $T\subseteq [m]$ with $|T|\le d-1$, and any Boolean test $f$ depending only on $b_T$, there exist parity-cut distributions $\mu_0,\mu_1$ on $\mathbb{F}_2^m$ with $\|\mu_0-\mu_1\|_{\mathrm{TV}}=1$ such that
\[
\mathbb{P}_{\mu_0}[f=1]=\mathbb{P}_{\mu_1}[f=1].
\]
\end{corollary}

\subsection{Conditional EntropyDepth lower bound template}

\begin{definition}[Clause-local transcript process]
A transcript process $(S_t)_{t\ge 0}$ is \emph{$q$-clause-local} if at each step $t$, the update from $S_t$ to $S_{t+1}$ reads at most $q$ new clause-side coordinates of $b\in\mathbb{F}_2^m$.
\end{definition}

For an adaptive process, the fixed-support corollary above does not by itself identify a single deterministic coordinate set supporting the entire transcript.  The required adaptive decision-tree/measurability statement is therefore separated explicitly.

\begin{definition}[Cumulative clause-dual support wall]
Fix a target success predicate for transcripts.  Say the process satisfies the \emph{clause-dual support wall at distance $d$} if, for every depth $t$, whenever the cumulative information exposed by the transcript is supported on fewer than $d$ clause coordinates, no positive-bias decision between the two parity-cut classes can be made from $S_t$.
\end{definition}

\begin{theorem}[Conditional: depth from clause-dual distance]
\textbf{Conditional.}
Let $A\in\mathbb F_2^{m\times n}$ have clause-dual distance
\[
d=d^\perp_{\mathrm{cl}}(A).
\]
Let $(S_t)_{t\ge0}$ be a $q$-clause-local transcript process, and assume the cumulative clause-dual support wall has been proved for the actual transcript semantics.
If a positive-bias correct decision is first possible from $S_t$, then
\[
d\le (t+1)q.
\]
Consequently, for $q>0$,
\[
t+1\ge \left\lceil\frac d q\right\rceil.
\]
In particular, if $d=\Omega(n)$ and $q=O(1)$, this conditional reduction yields $t=\Omega(n)$.
\end{theorem}

\begin{proof}
Through depth $t$, at most $(t+1)q$ fresh clause coordinates can have been exposed.
If $(t+1)q<d$, the assumed cumulative clause-dual support wall says that no positive-bias correct decision can be made from $S_t$, contradicting success at depth $t$.
Therefore $d\le (t+1)q$.
For $q>0$, the ceiling-form lower bound is the standard arithmetic reformulation.
\end{proof}

\paragraph{Formal boundary.}
The arithmetic implication
\[
\text{success at depth }t\;\wedge\;\text{clause-dual support wall}
\Longrightarrow d\le (t+1)q
\]
is mirrored by the Lean frontier theorem
\texttt{success\_forces\_cumulative\_support\_reaches\_dual\_distance}.
The remaining semantic obligation is to prove the cumulative clause-dual support wall for the actual adaptive transcript model.  No claim is made that $q<d$ alone prevents success for all time.
